% The abstract uses the conference abstract environment rather than a numbered section.
\begin{abstract}
Synthetic user prompts support automated evaluation and model orchestration, but apparent authorship could affect how a responder behaves. We examine prompt style, declared source, elicited authorship judgments, and activation interventions in an exploratory controlled study. Two API responders receive 180 benchmark requests spanning harmful assistance, benign assistance, and persona-conditioned opinions; a separate local model supports representation analysis. A formatting control preserves every underlying request byte and the complete case-insensitive word sequence. For Gemini 3.1 Flash Lite, this manipulation increases AI-authorship classifications by 64.4 percentage points on harmful requests and 88.1 points on benign requests, while refusal changes by 1.7 and 0.0 points. None of six primary model-by-domain behavioral comparisons is significant after Holm correction; uncertainty generally prevents equivalence claims. Broader paraphrasing fails strict semantic controls: only 7 of 124 successful benchmark pairs survive. In local Qwen3-8B, an activation classifier reaches 0.990 AUROC, but surface features reach 0.970 and surface residualization reduces activation AUROC to 0.665. Fixed-text steering yields no robust behavioral change specific to the style direction. An exploratory unconditional opinion-alignment decrease does not survive correction. These findings separate strong elicited source-judgment sensitivity from an unestablished behavioral effect and leave the general authorship-mediated causal hypothesis unresolved.
\end{abstract}
